\documentclass{article}
\usepackage{iclr2027_conference,times}
\input{math_commands.tex}
\usepackage[hidelinks]{hyperref}
\usepackage{url}
\usepackage{amsmath,amssymb,amsthm}

\newtheorem{theorem}{Theorem}
\newtheorem{lemma}{Lemma}
\newtheorem{corollary}{Corollary}
\newcommand{\Expo}{E}
\newcommand{\Budget}{B}

\title{Training-Path Dependence in Bounded Convex Dynamics:\\
Exposure Products, Clean-Tail Return, and an Affine Recursion}
\author{Anonymous}

\begin{document}
\maketitle

\begin{abstract}
This paper records three narrowly scoped ways in which a training-path
intervention can be analyzed. A fixed quadratic has an exact matrix-product
solution whose first-order log contraction depends on cumulative
learning-rate exposure, with a finite-step squared-step remainder. A
twice-differentiable, strongly convex clean objective contracts full
parameter error after an intervention is removed; a faster directional rate
requires an additional invariant-subspace condition. An affine scalar
conditional-mean recursion has an exact product law and an exponential upper
bound. These statements do not imply corresponding laws for logistic
regression, shuffled SGD, CNNs, nonconvex optimization, or real label noise.
The current public package makes no empirical or deployment claim.
\end{abstract}

\section{Introduction}
\label{sec:intro}

Training-data attribution and unlearning study relationships between a training
set and a trained model, including influence functions
\citep{koh2017understanding}, datamodels \citep{ilyas2022datamodels},
TRAK \citep{park2023trak}, and retraining-based unlearning
\citep{bourtoule2021sisa}. This paper addresses a more limited mathematical
question: when an objective is changed along a training trajectory, which
contraction quantities can be stated exactly?

The answer depends sharply on the model. We prove three results, each with its
own assumptions. T1 is a fixed-quadratic product identity. T2 is a
strong-convex clean-tail contraction result. T3 is an affine scalar
conditional-mean identity. The results share product structure but are not
interchangeable: neither T1 nor T3 establishes a law for nonlinear SGD, and
T2 is not a machine-unlearning guarantee for arbitrary neural networks.

\paragraph{Scope.}
The frozen public baseline contained historical numerical material without
co-located public code, output, configuration, seed, and preregistration
artifacts. That material is not part of this submission source. A separate,
uncompiled archive points to the immutable checkpoint containing it. A planned
two-scenario experiment is a falsification path only, not evidence.

\section{Formal setup}
\label{sec:setup}

Let $A$ denote a clean objective and let $D$ denote an intervention block.
The results below use only the objective active at the relevant step; they do
not assume that $D$ is a particular data type, or that an observed loss signal
identifies it. Write $\eta_t\geq0$ for a deterministic step size and
$E_{s:T}=\sum_{t=s}^{T-1}\eta_t$ for tail exposure. We reserve
the word return for distance to the minimizer of the clean objective, not for
accuracy, a representation, or a deployment outcome.

\section{Three bounded results}
\label{sec:theory}

\subsection{T1: fixed-quadratic exposure product}
\label{sec:t1}

\begin{lemma}[T1]
\label{thm:t1}
Let
\[
 Q(w)=\tfrac12(w-w^\star)^\top H(w-w^\star),\qquad
 0\prec\mu I\preceq H\preceq LI,
\]
and run $w_{t+1}=w_t-\eta_t\nabla Q(w_t)$ with
$0\leq\eta_tL<1$. Then
\begin{equation}
 w_T-w^\star=\prod_{t=0}^{T-1}(I-\eta_tH)(w_0-w^\star).
\label{eq:t1product}
\end{equation}
For an eigenvalue $\lambda_j$ of $H$, set
$P_j=\prod_t(1-\eta_t\lambda_j)$,
$E=\sum_t\eta_t$, and $q_j=\max_t\eta_t\lambda_j<1$. Then
\begin{equation}
 0\leq-\log P_j-\lambda_jE
 \leq\frac{\lambda_j^2\sum_t\eta_t^2}{2(1-q_j)}.
\label{eq:t1log}
\end{equation}
\end{lemma}

\begin{proof}
Since $\nabla Q(w)=H(w-w^\star)$, one step gives
$w_{t+1}-w^\star=(I-\eta_tH)(w_t-w^\star)$, and iteration proves
\eqref{eq:t1product}. Diagonalize the symmetric positive-definite $H$.
On an eigenvector with eigenvalue $\lambda_j$, the multiplier is $P_j$.
For $x\in[0,q_j]$,
\[
 0\leq-\log(1-x)-x
   =\sum_{k\geq2}\frac{x^k}{k}
 \leq\frac{x^2}{2(1-q_j)}.
\]
Apply this inequality to $x=\eta_t\lambda_j$ and sum over $t$.
This proves \eqref{eq:t1log}.
\end{proof}

T1 says that equal exposure yields the same first-order log contraction in
this fixed quadratic. It does not say that arbitrary finite schedules with
equal exposure have identical terminal iterates: the squared-step term
quantifies the permitted discrepancy. Replacing a step $\eta$ by $K$ steps of
size $\eta/K$ preserves exposure, whereas repeating $K$ original steps changes
exposure and is not a finite-budget invariance.

\subsection{T2: strong-convex clean-tail return}
\label{sec:t2}

\begin{theorem}[T2]
\label{thm:t2}
Let $L_A$ be twice continuously differentiable, $\mu$-strongly convex, and
$L$-smooth on every segment joining a tail iterate to its unique minimizer
$w_A^\star$. If $0\leq\eta_t\leq1/L$, then after switching to gradient
descent on $L_A$ at step $s$,
\begin{equation}
 \lVert w_T-w_A^\star\rVert_2
 \leq\prod_{t=s}^{T-1}(1-\eta_t\mu)
      \lVert w_s-w_A^\star\rVert_2
 \leq e^{-\mu E_{s:T}}\lVert w_s-w_A^\star\rVert_2.
\label{eq:t2full}
\end{equation}
If, in addition, $e\ne0$ and $\operatorname{span}(e)$ is invariant for every segment
average Hessian in the proof and its eigenvalue on $e$ is at least
$\lambda_e$, then
\begin{equation}
 |e^\top(w_T-w_A^\star)|
 \leq e^{-\lambda_eE_{s:T}}|e^\top(w_s-w_A^\star)|.
\label{eq:t2direction}
\end{equation}
\end{theorem}

\begin{proof}
For each $t\geq s$, the fundamental theorem of calculus gives
\[
 \nabla L_A(w_t)-\nabla L_A(w_A^\star)
 =\bar H_t(w_t-w_A^\star),\qquad
 \bar H_t=\int_0^1\nabla^2L_A(w_A^\star+
        r(w_t-w_A^\star))\,dr .
\]
The assumed Hessian bounds imply
$\mu I\preceq\bar H_t\preceq LI$. Therefore
\[
 w_{t+1}-w_A^\star=(I-\eta_t\bar H_t)(w_t-w_A^\star).
\]
All eigenvalues of $I-\eta_t\bar H_t$ lie in
$[0,1-\eta_t\mu]$ because $\eta_t\leq1/L$. Hence its spectral norm is at most
$1-\eta_t\mu$. Multiply these bounds from $s$ through $T-1$, then use
$1-x\leq e^{-x}$ to obtain \eqref{eq:t2full}.

For the directional statement, the added invariance premise makes $e$ an
eigenvector of every $\bar H_t$. Its scalar multiplier is at most
$1-\eta_t\lambda_e$; the same multiplication and exponential bound prove
\eqref{eq:t2direction}. A Rayleigh-quotient lower bound alone does not imply
this shared invariant-subspace premise.
\end{proof}

Thus T2 is a full-parameter result under explicit strong-convexity conditions.
Its directional refinement is conditional and must not be inferred from
observations in another model class.

\subsection{T3: affine scalar conditional-mean product}
\label{sec:t3}

\begin{theorem}[T3]
\label{thm:t3}
Let $(u_t)$ be an integrable scalar process adapted to
$(\mathcal F_t)$. Fix $\lambda>0$ and deterministic step sizes satisfying
$0<\eta_t\lambda<1$. Suppose
\begin{equation}
 u_{t+1}=(1-\eta_t\lambda)u_t+\eta_t\xi_t,\qquad
 \mathbb E[\xi_t\mid\mathcal F_t]=0.
\label{eq:t3affine}
\end{equation}
Then for any $s<T$,
\begin{equation}
 \mathbb E[u_T]=\mathbb E[u_s]\prod_{t=s}^{T-1}
 (1-\eta_t\lambda),\qquad
 |\mathbb E[u_T]|
 \leq|\mathbb E[u_s]|e^{-\lambda E_{s:T}}.
\label{eq:t3product}
\end{equation}
\end{theorem}

\begin{proof}
Taking conditional expectations in \eqref{eq:t3affine} and then total
expectations yields
\[
 \mathbb E[u_{t+1}]=(1-\eta_t\lambda)\mathbb E[u_t].
\]
Repeated substitution gives the product identity in \eqref{eq:t3product}.
Every factor is positive and $1-x\leq e^{-x}$ for $x\geq0$, so multiplying
the inequalities gives the exponential upper bound.
\end{proof}

T3 is exact for the displayed scalar recursion. It does not assert that
logistic-SGD, random reshuffling, or a nonlinear accuracy functional obeys
\eqref{eq:t3affine}.

\section{Empirical status and falsification path}
\label{sec:empirical-status}

No current empirical claim is made. PLAN-ENABLE-001 specifies a future
falsification experiment for a training-visible enable policy in two real
scenarios: natural human-label noise and a clean benign-hard cohort. It must
freeze code, data and split hashes, comparator arms, seeds, environment,
event schema, and forbidden-field checks before execution. Success on one
scenario alone would not establish a deployment claim; failure is an
informative boundary for that particular policy, not a refutation of T1--T3.

\section{Related work}
\label{sec:related}

Influence functions \citep{koh2017understanding}, datamodels
\citep{ilyas2022datamodels}, and TRAK \citep{park2023trak} provide relevant
ways to relate training data to learned-model behavior. SISA
\citep{bourtoule2021sisa} is a different retraining-based unlearning
construction. Work on SGD noise and implicit regularization
\citep{mandt2017sgd,smith2017bayesian,smith2021origin} motivates stating the
conditional-mean premise in T3 rather than treating it as an SGD identity.
Loss-mixture and noisy-label work \citep{arazo2019unsupervised,jiang2018mentornet,
han2018coteaching,wei2022cifar10n} motivates the planned experiment, not a
present policy result. The closest-work matrix records the checked metadata
and the deliberately narrow relationship claims.

\section{Limitations}
\label{sec:limitations}

The scope is mathematical rather than empirical. T1 is fixed quadratic. T2
requires smooth strong convexity along the relevant segments, a stable
step-size bound, and a separate invariant-subspace condition for its
directional refinement. T3 requires the affine scalar recursion and
conditional-zero-mean innovation. None of these assumptions automatically
holds for logistic training, shuffled SGD, neural networks, noisy labels, or
data-selection systems. The public package also lacks a same-version
experiment evidence chain; remote candidate results are not attributed to
this manuscript.

\section{Conclusion}

The three results establish exact or bounded product relations only within
their stated regimes. Their shared form motivates careful future testing of
path interventions, but it does not license a broader theory or a deployment
recommendation. The next empirical step is a registered, matched-path
falsification test, not recovery of unsupported historical numbers.

\section*{Reproducibility Statement}

This self-contained source package includes its TeX dependencies and
bibliography and is built in isolation. Formal assumptions, statements, and
proofs are active in this source. The separate execution-readiness audit names
the artifacts required before PLAN-ENABLE-001 can produce evidence; no
unexecuted protocol or remote-only result is reported as an observation.

\section*{Ethics Statement}

This revision reports no newly run human-subject or dataset study. Any future
execution must verify dataset licenses, versions, and use conditions, and the
policy must not access clean, protected, test, or oracle labels at decision
time.

% AI Use Statement — ICLR 2027 submission.  This statement is deliberately
% revision-specific and supersedes inherited agent/workflow descriptions.

\section*{AI Use Statement}
\label{sec:ai-use}

This manuscript and its audit package were prepared with substantial assistance
from LLM-based coding agents, including Codex and earlier multi-agent research
workflows, under human supervision. The use of such systems included drafting,
revision, repository inspection, provenance and citation checks, build checks,
and internal review planning. No LLM-generated assertion was treated as
evidence without an identified source artifact or an explicit ``planned'' label.

\paragraph{Tasks performed with AI assistance.}
\begin{itemize}
  \item Drafting and revising prose, equations, figures, tables, and this
  disclosure; generating alternative paper stories and revision plans.
  \item Inspecting the public baseline and a separately frozen remote candidate;
  computing hashes and diffs; checking source-package dependencies; and running
  safe, read-only consistency/configuration verifiers.
  \item Searching for and cross-checking primary or official bibliographic
  metadata, then building a claim--evidence map and closest-work matrix.
  \item Designing (but not submitting or executing) the planned matched-path
  experiment and auditing its data, environment, launcher, comparator, seed,
  and output-schema readiness.
  \item Conducting or assisting internal manuscript audits. Such internal
  reviews are not peer review and do not establish experimental results,
  theorem correctness, or acceptance prospects.
\end{itemize}

\paragraph{Human responsibility and controls.}
\begin{itemize}
  \item Human supervisors specify the task, authorization boundaries, resource
  limits, and submission decisions. This revision did not authorize a GPU,
  remote, or formal experimental run.
  \item Human authors must verify all final claims, citations, proofs, data-use
  permissions, and venue compliance, and take full responsibility for the
  submission and any remaining errors.
\end{itemize}

The current package intentionally distinguishes the public manuscript version
from a different remote candidate. Remote scripts and results are not used as
evidence for public-version claims unless a claim-level reconciliation closes
that gap. Accordingly, this disclosure does not imply that any historical
agent-produced numerical claim has been independently reproduced.


\bibliographystyle{iclr2027_conference}
\bibliography{references}

\appendix
\section{Proof-scope checklist}

T1 uses a fixed symmetric positive-definite Hessian and its eigendecomposition.
T2 uses segment Hessians bounded between $\mu I$ and $LI$; its directional
line is conditional on simultaneous invariance. T3 takes conditional
expectations of the stated scalar recursion. These are the only proof
dependencies used by the active manuscript. Historical figures, tables,
scripts, and numerical claims are not imported, hidden, or compiled here.

\end{document}
